\subsection*{§. 200.\\ 零曲率直纹面与 Clifford 曲面。}

欧几里得几何中的平行性与椭圆几何中的平行性之间，还有一种显著的类比，表现于以下性质。任取一个 Clifford 直线汇中的一重无穷多条直线，亦即考察一个直纹面 \(S\)，其母线在 Clifford 意义下彼此平行。由该直线汇确定的连续滑移使 \(S\) 沿自身滑动，而每条母线各自沿自身滑动。因此，把 \(S\) 的母线取作 \(u=\) 常数，把正交轨线取作 \(v=\) 常数，则由 §. 104（第 231 页）的定理，其线元必有形式
\[
 ds^2=du^2+\varphi(u)dv^2.
\]
但因曲线 \(v=\) 常数为测地线，必有 \(\varphi(u)=\) 常数，从而 \(ds^2\) 具有普通欧几里得平面的形式\srcfn{1}{正文中的论证也可表述如下。任取母线的一条正交轨线 \(c\)，滑移会把它移到其余所有正交轨线上。因此，母线在这些正交轨线 \((c)\) 上截得的弧相等，故而（参见 §. 89 的注）曲线 \((c)\) 本身都是测地线。于是，曲面 \(S\) 具有两族彼此正交的测地线，故其曲率为零。}。于是得到定理：\emph{椭圆空间中包含于一个 Clifford 直线汇内的每个直纹面，都是总曲率为零的曲面。} 这个定理把 Clifford 直纹面与欧几里得空间中的柱面对应起来。

特别考察一个 Clifford 直线汇中的直纹面，其母线 \(g\) 从同一条不属于该直线汇的直线 \(r\) 的各点发出。它们组成一个特殊直纹面；其上存在两族不同的、各自彼此平行的母线 \((g)(r)\)，第一族为右旋平行，第二族为左旋平行（参见 §. 196 末尾）。在由两对母线 \((g_1,g_2),(r_1,r_2)\) 构成的每个四边形中，相对边相等，相对角也相等（在椭圆度量下）。这个曲面称为 \emph{Clifford 曲面}；显然它有两个不同种类、彼此可交换且使曲面变为自身的滑移。Clifford 正是由此推出该曲面的曲率为零；不过如上所见，这个性质其实属于每个具有平行母线的直纹面。

下列讨论还给出 Clifford 曲面的另一种显著生成方式\srcfn{1}{参见我的论文 \emph{Sulle superficie a curvatura nulla in geometria ellittica}（《Annali di matem.》，1895）。}。Clifford 曲面所具有的两个滑移确定两条关于绝对形互为极线的直线，它们既平行于第一族母线，也平行于第二族母线。以 \(r,r'\) 表示这两条直线，称为 Clifford 曲面的轴；在两个滑移中的任一个作用下，它们都沿自身滑动，因而在由这两个滑移复合成的任意运动下也沿自身滑动。现在把沿 \(r\) 的一个右旋滑移与角幅相等而相反的一个左旋滑移复合；所有通过 \(r\) 的平面都将沿自身滑动，而它们在 Clifford 曲面上截出的截线也随之如此，因为该曲面在所考察运动下沿自身滑动。可是，在椭圆平面中一个使某直线沿自身滑动的运动下，其余各点的轨迹都是圆；这些圆与该直线测地平行，圆心为该直线的极点。因此，通过轴 \(r\) 的平面（从而垂直于第二轴 \(r'\)）在 Clifford 曲面上截出的曲线，是半径相等的圆，其圆心分布在第二轴 \(r'\) 上。显然，绕 \(r\) 的纯\emph{旋转}使这些圆彼此交换。因此，可以如下生成 Clifford 曲面：\emph{任取一个圆，并在其平面内作圆心的极线；令该平面绕这条直线旋转，所取圆便描出 Clifford 曲面。}

显然，Clifford 曲面还有第二族圆；并且容易证明，这两族圆恰是该曲面的曲率线，因而是两族母线所成双重系统的角平分线。

